\section{Simplicial homotopy theory and group actions}

This section records results on ordinary simplicial sets, simplicial categories,
and group actions used in the formalization of Theorems A--D.
Write $C_2=\{1,\omega\}$, where $\omega^2=1$.
Two distinct realization functors occur here:
$\operatorname{Real}_{Q_{\mathrm{rc}}}:\sSet\to\sSet$ is built from the
right-cone cosimplicial object $Q_{\mathrm{rc}}$, whereas
$|-|:\sSet\to\mathrm{TopCat}$ denotes topological realization.
The symbol $Q_{\mathrm{rc}}$ names an auxiliary object in the Lean formalization;
it is different from the paper's self-conjugate simplices $Q^n$.

\begin{resultblock}
\begin{theorem}[A Kan--Quillen self-equivalence from $Q_{\mathrm{rc}}$]
In every universe $u$, the left Kan extension
$\operatorname{Real}_{Q_{\mathrm{rc}}}:\sSet\to\sSet$ of the original
right-cone cosimplicial object $Q_{\mathrm{rc}}$ and its right adjoint
$\operatorname{Sing}_{Q_{\mathrm{rc}}}$ form a Quillen equivalence for the
ordinary Kan model structure.
In particular, for every Kan complex $X$, the counit
$\operatorname{Real}_{Q_{\mathrm{rc}}}\operatorname{Sing}_{Q_{\mathrm{rc}}}X\to X$
is a weak equivalence.
This is not the adjunction between topological realization and the ordinary
singular complex of a topological space.
\end{theorem}
\Evidence{DaggerActualQQuillenEquivalence.isQuillenEquivalence; DaggerActualQKanCounitWeakEquivalence.counit_weakEquivalence}{ActualQQuillenEquivalence.lean:17; ActualQKanCounitWeakEquivalence.lean:24}
\Origin{The construction and comparison for the original right-cone object $Q_{\mathrm{rc}}$ are proved in this formalization, using the existing Kan model structure and Mathlib's Quillen-equivalence criterion. This is an ordinary simplicial-set result, distinct from the dagger rigidification result in Theorem B.}
\end{resultblock}

\begin{resultblock}
\begin{theorem}[The projective model structure on simplicial $C_2$-actions]
In every universe $u$, the category
$\Fun(BC_2,\sSet)_{\proj}$ of simplicial $C_2$-actions admits a projective
model structure.
If $U$ is the forgetful functor and $f$ is an equivariant map, then
\begin{align}
 f\text{ is a weak equivalence}&\iff U(f)\text{ is a Kan weak equivalence},\\
 f\text{ is a fibration}&\iff U(f)\text{ is a Kan fibration}.
\end{align}
Its cofibrations form the left lifting class generated by the free simplicial
boundary inclusions.
The assertion concerns this $C_2$-action category; it does not claim the same
construction unconditionally for every group.
\end{theorem}
\Evidence{DaggerOriginalSimplicialActionModelCategory.modelStructure; DaggerOriginalSimplicialActionModelCategory.weakEquivalence_iff_underlying; DaggerOriginalSimplicialActionModelCategory.fibration_iff_underlying}{OriginalSimplicialActionModelCategory.lean:126,58,131}
\Origin{The free-boundary and free-horn generators, finite-source smallness, and closure of weak equivalences are verified for this action category. The construction uses Mathlib, the ported TopPackage machinery, and the ordinary Kan model structure; it does not present that general machinery as a new theorem.}
\end{resultblock}

\begin{resultblock}
\begin{theorem}[The Borel model structure on topological $C_2$-actions]
In every universe $u$, the category of $C_2$-actions in
$\mathrm{TopCat}$ admits a Borel model structure.
An equivariant map is a weak equivalence, respectively a fibration, exactly
when its underlying map has that property in the existing
$\mathrm{TopCat}$ model structure.
The cofibrations coincide with the original Borel left lifting class.
The result uses that particular topological model structure and makes no claim
about a different category of topological spaces.
\end{theorem}
\Evidence{DaggerOriginalBorelModelCategory.modelStructure; DaggerOriginalBorelModelCategory.weakEquivalence_iff_underlying; DaggerOriginalBorelModelCategory.fibration_iff_underlying; DaggerOriginalBorelModelCategory.cofibration_iff_original}{OriginalBorelModelCategory.lean:73,79,85,91}
\Origin{This is a concrete model construction for the action category, using free realized boundaries and horns and finite free realization sources. The underlying topological q-model and package machinery are reused from pinned TopCatModelCategory sources; their upstream proofs are not claimed as new independent proofs here.}
\end{resultblock}

\begin{resultblock}
\begin{theorem}[Cofibrancy of $E_{\bullet}C_2$ and contractibility of $EC_2$]
In every universe $u$, the $C_2$-action on the simplicial set
$E_{\bullet}C_2$ is cofibrant in the projective model structure above.
Its topological realization $EC_2:=|E_{\bullet}C_2|$ is contractible as an
\emph{underlying} topological space.
Thus the two conclusions are $\operatorname{IsCofibrant}(E_{\bullet}C_2)$ and
$\operatorname{ContractibleSpace}(EC_2)$.
The first follows from an equivariant relative cell complex built from free
boundary cells; the second follows by realizing an explicit contraction.
The latter does not assert equivariant contractibility or contractibility of
fixed-point spaces.
\end{theorem}
\Evidence{DaggerOriginalEC2SimplicialProjectiveCofibrancy.originalSimplicialEC2_isCofibrant; DaggerOriginalEC2RealizationContractibility.realizedContractibleSpace; DaggerOriginalEC2EquivariantRelativeCellComplex.equivariantRelativeCellComplex}{OriginalEC2SimplicialProjectiveCofibrancy.lean:60; OriginalEC2RealizationContractibility.lean:58; OriginalEC2EquivariantRelativeCellComplex.lean:84}
\Origin{The action, skeletal pushouts, equivariant relative cell complex, and simplicial contraction are constructed for the original $E_{\bullet}C_2$. General cell-complex and topological-realization APIs are reused from Mathlib and TopCat.}
\end{resultblock}

\begin{resultblock}
\begin{proposition}[Kan fibers of prefix mapping spaces]
Let $X\in\sSet$ be a quasi-category in any universe $u$.
For every $k\geq0$, every $k$-simplex $\sigma$ of $X$, and every vertex $b$,
the fiber $\operatorname{prefixMapping}(X,\sigma,b)$ of the original coslice
projection over $b$ is a Kan complex.
The case $k=0$ includes the Kan property of the ordinary left mapping space.
The stated result retains the quasi-category hypothesis on $X$.
\end{proposition}
\Evidence{DaggerActualQuasicategoryMappingFibres.prefixMapping_kan; DaggerActualQuasicategoryMappingFibres.leftMapping_kan}{ActualQuasicategoryMappingFibres.lean:32,37}
\Origin{The proof uses lifting for the original coslice projection and the Kan criterion for its fibers. The statement covers every prefix simplex and terminal vertex, rather than an individual simplicial face calculation.}
\end{resultblock}

\begin{resultblock}
\begin{theorem}[Weak-equivalence invariance of right mapping spaces]
Let $X,Y$ be quasi-categories in any universe $u$, let $f:X\to Y$ be a
categorical weak equivalence, and let $a,b\in X_0$.
The induced map of the original right-cone mapping spaces
\begin{align}
 \Map^{R}_{X}(a,b)
  &\longrightarrow \Map^{R}_{Y}(f(a),f(b))
\end{align}
is a weak equivalence of ordinary simplicial sets.
No monomorphism or Kan-fibration assumption is imposed on $f$.
\end{theorem}
\Evidence{DaggerActualRightMappingWeakInvariance.map_weak_of_categorical_weak}{ActualRightMappingWeakInvariance.lean:36}
\Origin{The proof uses the original right-cone mapping spaces, a span for categorical weak equivalences, and lifting. It is distinct from the unitary-core criterion in the dagger Theorem C.}
\end{resultblock}

\begin{resultblock}
\begin{theorem}[Right-cone mapping spaces and $Q_{\mathrm{rc}}$-singular enriched Homs]
Let $D\in\sCat$ be an ordinary simplicial category in any universe $u$,
and let $v,w$ be vertices of its coherent nerve.
Write $\bar v,\bar w$ for their corresponding objects of $D$.
There is an isomorphism of simplicial sets
\begin{align}
 \operatorname{rightConeSimplices}_{D}(v,w)
  &\cong \operatorname{Sing}_{Q_{\mathrm{rc}}}(D(\bar v,\bar w)).
\end{align}
The comparison is bijective in every simplicial degree and respects the
simplicial operators.
No additional naturality statement in $D,v,w$ is asserted here.
\end{theorem}
\Evidence{DaggerActualRightConeSingularComparisonIso.comparisonIso; DaggerActualRightConeSingularComparisonIso.comparison_app_bijective}{ActualRightConeSingularComparisonIso.lean:43,17}
\Origin{The isomorphism compares the original cone evaluation with the original $Q_{\mathrm{rc}}$-singular construction. It concerns mapping objects and is separate from the ordinary rigidification--coherent-nerve adjunction itself.}
\end{resultblock}
